\documentclass[11pt,a4paper]{article}
\usepackage[margin=2cm]{geometry}
\usepackage{amsmath,amssymb}
\setlength{\parskip}{0.5em}
\setlength{\parindent}{0pt}

\begin{document}

\begin{center}
  {\LARGE \bfseries Angles in the Same Segment --- Answer Key}\\[0.25em]
  \small Angles subtended by the same chord at the circumference, in the same segment, are equal.
\end{center}

\hrule
\vspace{1em}

\section*{1. Identifying the segments}

A \textbf{segment} is the region bounded by a chord and one of the arcs it cuts off.  
A \textbf{sector} is the region bounded by two radii and the arc between them.

\begin{center}
\begin{tabular}{c|c|c|c}
 & Left & Middle & Right \\ \hline
Row 1 & segment & sector & neither \\
Row 2 & segment & sector & segment \\
Row 3 & segment & sector & neither \\
Row 4 & segment & sector & segment
\end{tabular}
\end{center}

The shaded region in Row 2, Right is the \emph{major} segment; it is still a segment.

\section*{2. Angles in a given segment}

Place the marked angle vertices on the shaded arc. Any points on that same shaded arc give equal angles standing on chord \(AB\).

\begin{itemize}
  \item[(a)] One angle: put vertex \(C\) on the shaded arc.
  \item[(b)] Two angles: put vertices \(C\) and \(D\) on the shaded arc.
  \item[(c)] Three angles: put vertices \(C\), \(D\), and \(E\) on the shaded arc.
  \item[(d)] Ten angles: put vertices \(P_1,\dots,P_{10}\) on the shaded arc.
\end{itemize}

In (a) and (b) the shaded segment is the minor segment.  
In (c) and (d) the shaded segment is the major segment.  
In every case, the marked angles are equal because they are in the same segment.

\section*{3. Interpreting the theorem}

In each pair, shade the segment that contains the arc on which the angle vertices \(C\) and \(D\) lie.

\begin{itemize}
  \item[(a)] \(C\) and \(D\) lie on the major arc, so shade the \textbf{major segment}.
  \item[(b)] \(C\) and \(D\) lie on the major arc, so shade the \textbf{major segment}.
  \item[(c)] \(C\) and \(D\) lie on the minor arc, so shade the \textbf{minor segment}.
  \item[(d)] \(C\) and \(D\) lie on the major arc, so shade the \textbf{major segment}.
\end{itemize}

\end{document}